% MergeMap: A Network Analysis of Model Merging on Hugging Face
% COMP 6917 Complex Networks, Memorial University of Newfoundland, Fall 2026
% Compile from this folder:  pdflatex MergeMap_Report.tex  (run twice for references)
% Figures are read from ../results/, so keep the project folder layout unchanged.

\documentclass[conference]{IEEEtran}
\IEEEoverridecommandlockouts

\usepackage[T1]{fontenc}
\usepackage{amsmath,amssymb}
\usepackage{graphicx}
\usepackage{booktabs}
\usepackage{array}
\usepackage{url}
\usepackage[hidelinks]{hyperref}
\usepackage{cite}

\graphicspath{{../results/}}
\newcommand{\model}[1]{\texttt{\small #1}}

\begin{document}

\title{MergeMap: A Network Analysis of\\Model Merging on Hugging Face}

\author{\normalsize Md Minhajul Abedin \texttt{mminhajula@mun.ca}\\
COMP 6917 Complex Networks \quad $|$ \quad December 07, 2026}

\maketitle

\begin{abstract}
Hugging Face lets every model declare the models it was built from. A merged
model has two or more such parents, so its lineage is a graph and not a tree.
We build the lineage network of 942,088 models and 929,941 ``derived from''
links from a pinned snapshot of the public Hub metadata, and study it with the
methods of the course: components, degree distributions, cocitation, PageRank,
betweenness, breadth-first search, modularity, random-graph comparisons and
network formation models. Merges join the network: the largest weakly
connected component grows 9.6 times, from 43,660 to 418,549 models, when merge
links are added. The number of children per model follows a power law with
exponent 1.92, and 94 models (the top 0.01\%) hold 42.1\% of all links. Merges
almost never mix families (2.9\% of the merges whose family is known), and 48.7\%
of merges combine parents that already share an ancestor. Shuffled networks
that keep every degree reach none of these values in 100 runs. Merges are only
1.6\% of the models but carry 82.6\% of the lineage betweenness. Finally, under a
full-inheritance assumption, merges add 31.7\% to the number of models that a
flaw in a base model could reach, and the top 100 models by PageRank cover
52.1\% of all exposure, against 0.2\% for 100 random models with a child.
\end{abstract}

\begin{IEEEkeywords}
complex networks, model merging, Hugging Face, lineage graph, power law,
PageRank, betweenness, modularity, preferential attachment
\end{IEEEkeywords}

% ---------------------------------------------------------------------------
\section{Introduction}

Open model hubs work like supply chains. A team fine-tunes a base model,
another team quantizes it, a third merges it with two other models, and the
result is fine-tuned again. Model merging combines the weights of several
trained models without further training. It is cheap and popular, so many
recent models have several parents. Each Hugging Face model card can declare
its parents and the relation to them (fine-tune, adapter, quantization or
merge). Laufer et al.\ \cite{laufer2025} studied these declarations as family
trees. A merge breaks the tree: it has more than one parent, and it can connect
lineages that were separate. This project asks what merges do to the network
as a whole.

We follow the six-step workflow of the course: data set, graph model, metrics,
algorithms, knowledge and understanding. The project answers four research
questions.

\textbf{RQ1.} How much do merges connect separate model families?

\textbf{RQ2.} Who merges with whom, and how often do close relatives merge?

\textbf{RQ3.} Which models act as bridges or hubs?

\textbf{RQ4.} How many models inherit a flawed base model, and how much do
merges add?

We also added a short network-formation study (Topic~5-5): does the network
grow by preferential attachment, and do merges copy earlier merges?

All metrics and algorithms come from the lecture topics, with two exceptions
that the submitted proposal promised: family assortativity and
degree-preserving random graphs. Both are taken from Newman \cite{newman2010}.
Table~\ref{tab:methods} lists every method and its source.

% ---------------------------------------------------------------------------
\section{Related Work}

Laufer, Oderinwale and Kleinberg \cite{laufer2025} built the family trees of
Hugging Face models from the same \texttt{baseModels} field and studied how
traits change from parent to child. They treat lineage as trees and do not
study merges as network links. Horwitz et al.\ \cite{horwitz2025} propose charting
the Hub as a model atlas: a directed graph of models linked by the weight
transformations between them. Wei et al.\
\cite{wei2026} study which models are disruptive in the open-source LLM
community. Our work differs in two ways: it treats merges as the links that
join lineages, and it uses classic network baselines (random graphs and
degree-preserving shuffles) to test whether the patterns are more than an
effect of the degrees. The network models we use are those of Price
\cite{price1976}, Barab\'asi and Albert \cite{barabasi1999} and the textbooks
\cite{barabasi2016,newman2010}.

% ---------------------------------------------------------------------------
\section{Step 1: Data Set}

The source is the \texttt{models.parquet} file of the public data set
\texttt{cfahlgren1/hub-stats} (Apache-2.0), which holds the metadata of every
public model on the Hub. We pinned one snapshot: commit
\texttt{1a8f69d}, uploaded on 22~October~2026, 1.64~GB, with its SHA-256 hash
stored in the project. Each row has the model id, author, creation time and
the \texttt{baseModels} field, which gives the relation and the list of
declared parents. There are four relations: \emph{fine-tune} (the weights are
trained further), \emph{adapter} (a small set of added weights, such as LoRA,
trained on a frozen base), \emph{quantized} (the same weights stored at lower
precision) and \emph{merge} (the weights of several parents are combined).
The Hub back-fills the creation time of all models made before 2~March~2022 to
that one date, so time-based results use only the models created after it.

Table~\ref{tab:cleaning} shows the cleaning steps. The file has 3,177,876 rows;
13 duplicate ids were removed. Of the 3,177,863 models, 902,519 declare at least
one parent. We removed self-loops and links to a parent created after the
child (these are mistakes or later re-uploads). We kept the 14,006 links to
3,431 parents that are no longer public, as ``ghost'' nodes, because removing
them would cut real lineages. No model declares two parents unless the relation
is ``merge''. Of the 15,310 models that declare a merge, 264 have fewer than two
parents in the clean graph; as planned in the proposal, they count for reach (RQ4)
but not for the mixing and inbreeding results, which use the 15,046 merges with
two or more parents.

\begin{table}[t]
\caption{Data cleaning of the declared parent links.}
\label{tab:cleaning}
\centering
\small
\begin{tabular}{@{}lr@{}}
\toprule
Step & Count \\
\midrule
Rows in the snapshot & 3,177,876 \\
Models after removing duplicate ids & 3,177,863 \\
Models that declare a parent & 902,519 \\
Raw child $\to$ parent links & 935,443 \\
\quad removed: self-loops & 2,072 \\
\quad removed: parent newer than child & 3,430 \\
\quad kept: links to 3,431 ghost parents & 14,006 \\
Clean links & \textbf{929,941} \\
\quad fine-tune / adapter & 358,180 / 332,026 \\
\quad quantized / merge & 191,705 / 48,030 \\
Nodes (models with a link, incl.\ ghosts) & 942,088 \\
Merge models (2+ parents) & 15,046 \\
\bottomrule
\end{tabular}
\end{table}

\textbf{Check against earlier work.} Laufer et al.\ \cite{laufer2025} count
191,072 fine-tune links in their snapshot of 1.86 million models. Cut at the
date when our snapshot reaches 1.86 million models (22~August~2025), our graph
has 191,155 (a 0.04\% difference; the cut is approximate because deleted models
are missing from our snapshot). We also find the same most-derived model
(\model{FLUX.1-dev}), and a maximum depth of 39 when each merge keeps only its
first-listed parent, against their ``nearly 40 generations'' under the same
rule. This tells us the pipeline reads the field correctly.

% ---------------------------------------------------------------------------
\section{Step 2: Graph Model}

\textbf{Lineage graph.} Each model is a node. Each declared parent gives a
directed edge \emph{child $\to$ parent}, read as ``is derived from''. With this
direction the in-degree of a model is its number of children, which is the
measure of influence used in RQ1 and RQ3. The proposal drew edges the other
way; we changed it so that in-degree, PageRank and the citation analogy of the
lectures (a child ``cites'' its parent) all point the same way. A child must be
newer than its parent, so the graph should be a directed acyclic graph (DAG).
We checked this with our own Kahn topological sort (Topic~2): every node was
ordered, so there is no cycle.

We compare two graphs: $G_{\text{full}}$ with all edges, and
$G_{\text{nomerge}}$ in which every merge model keeps no parent links. The
difference between the two is the effect of merges.

\textbf{Families.} A root is a model with no parent. We gave each root a family
from a fixed list of 41 name patterns (e.g.\ \texttt{llama}, \texttt{qwen},
\texttt{mistral}, \texttt{gemma}, \texttt{flux}), written before any result was
computed, and labelled the rest \emph{unknown}. A non-root model takes the
family of its roots, or \emph{mixed} if its roots belong to different named
families. 21,165 of the 44,867 roots get a named family, and 85.3\% of the
non-root models get a named family or \emph{mixed}. The remaining 14.7\%, and
23.5\% of all merge parents, are unknown, so every family result below is
computed on labelled models and stated as such. The largest labels among
non-root models are Qwen (249,015), Llama (104,203), Gemma (85,641), BERT
(54,887), FLUX (46,217) and Mistral (44,174); only 279 models are \emph{mixed}.

\textbf{Cocitation network.} For RQ2 and RQ3 we project the merges onto their
parents. Two models are linked if they are parents of the same merge, and the
weight is the number of merges they share. This is the cocitation matrix
$C = A^{\mathsf T}A$ restricted to merge rows (Topic~2). It has 10,125 nodes,
47,300 links, total weight 112,141 and mean degree 9.34. The heaviest link joins
\model{Mistral-7B-Instruct-v0.2} and \model{Mistral-7B-v0.1} (247 merges).

% ---------------------------------------------------------------------------
\section{Steps 3 and 4: Metrics and Algorithms}

Table~\ref{tab:methods} maps each question to its method and lecture topic.
We wrote our own code for the methods that are central to the questions
(topological sort, depth, PageRank, inbreeding search, copying, the attachment
rate and the two merge shuffles), and used python-igraph \cite{igraph} for
components, betweenness, clustering and the rewiring of the cocitation network. Communities come from the Leiden algorithm
\cite{traag2019} in the \texttt{leidenalg} package.

\begin{table}[t]
\caption{Metrics and algorithms, with their source.}
\label{tab:methods}
\centering
\footnotesize
\setlength{\tabcolsep}{3pt}
\begin{tabular}{@{}p{0.9cm}p{4.6cm}p{2.1cm}@{}}
\toprule
RQ & Method & Source \\
\midrule
1 & DAG check by topological sort; depth & Topic 2 \\
1 & Weakly connected components, giant component & Topics 2, 4 \\
1 & Children distribution: log binning, cumulative slope, top-heavy share & Topic 4 \\
2 & Cocitation network $A^{\mathsf T}A$ & Topic 2 \\
2 & Modularity $Q$ with families as groups & Topic 5-3 \\
2 & Family assortativity $r$ & Newman \cite{newman2010} \\
2 & Clustering, giant component, path length vs.\ $G(n,m)$ & Topic 5-4 \\
2 & Inbreeding search (BFS up to depth 3) & Topic 5 \\
3 & PageRank, solved on the DAG & Topic 3 \\
3 & Betweenness (exact, directed and on cocitation) & Topic 5 \\
3 & Community detection (Leiden) & Topic 5-3 \\
4 & Reach by BFS; hub vs.\ random removal & Topics 4, 5 \\
Ext. & Attachment rate $R(k)$, Price's model, copying & Topic 5-5 \\
All & Degree-preserving shuffles & Newman \cite{newman2010} \\
\bottomrule
\end{tabular}
\end{table}

\textbf{Power law.} Following Topic~4, the children distribution is shown with
logarithmic bins, and the exponent $\alpha$ is taken from the slope of the
cumulative distribution $P(K\ge k)\sim k^{-(\alpha-1)}$ on log-log axes. A
maximum-likelihood fit (\texttt{powerlaw} package) is reported only as a
cross-check.

\textbf{PageRank.} Eigenvector centrality is zero for every node of a DAG
(Topic~3), so we use PageRank $x = \alpha A D^{-1}x + \beta\mathbf{1}$ with
$\alpha=0.85$: a model passes 85\% of its score to its parents, split equally.
In a DAG this system can be solved exactly in one pass, by handling models from
the deepest to the shallowest, since a child is always deeper than its parents.
The result matches the matrix inverse $(I-\alpha AD^{-1})^{-1}\beta\mathbf{1}$
(our Assignment~3 code) to $1.1\times10^{-11}$ on a 986-node test piece, and is
proportional to igraph's PageRank to within $7.5\times10^{-13}$.

\textbf{Inbreeding.} Two parents of a merge are \emph{related within $d$} if
their sets of ancestors within $d$ generations (each including the model
itself) overlap. A breadth-first search up the graph (Topic~5) finds these sets.
A pair is of \emph{ancestor} type if one parent is an ancestor of the other, and
of \emph{common-ancestor} type otherwise (siblings, cousins). A merge is
inbred within $d$ if any pair of its parents is related within $d$. We fixed
this definition before computing any result.

\textbf{Assortativity.} For family labels, $Q=\sum_i(e_{ii}-a_i^2)$ is the
modularity of Topic~5-3 with families as groups, and the assortativity
coefficient is $r = Q/(1-\sum_i a_i^2)$ \cite[eq.~7.80]{newman2010}, so $r=1$
means that links only join models of the same family.

\textbf{Baselines.} A pattern only matters if it is stronger than what the
degrees alone produce. We use three degree-preserving shuffles \cite{newman2010},
each run 100 times with seed 6917 and ten swap attempts per edge:
(A) swap parents between merge links, keeping every merge's number of parents
and every parent's number of merge children, and rejecting swaps that create a
duplicate parent or a cycle; (A2) the same, but swaps only inside a family, so
that family membership is kept as well; (B) igraph's \texttt{rewire} on the
cocitation network. We report the baseline mean, standard deviation and
$z$-score.

\textbf{Network formation.} We replay all 929,941 links in the order the
children were created. For a model with $k$ children, the attachment rate
$R(k)$ is the number of new links that went to models with $k$ children,
divided by the number of such models at the time. Linear preferential
attachment predicts $R(k)\propto k$ (Topic~5-5). For copying, each merge is
compared with every earlier merge that shares at least two parents; it
\emph{copies} if the best Jaccard similarity of the two parent sets is at
least 0.5, and \emph{repeats} a recipe if the sets are equal.

\textbf{Reproducibility.} All row orders, ties and seeds are fixed, and the best
of 20 Leiden runs is kept. Two full runs from scratch gave identical output
files. A full run takes about 14 to 19 minutes on a laptop.

\textbf{Changes from the proposal.} (1)~Edges point child~$\to$~parent.
(2)~Families come from a fixed list with an \emph{unknown} label. (3)~Ghost
parents are kept. (4)~The power law is fitted on the number of children only.
(5)~Path lengths are measured on the cocitation network, where they have a
meaning, and lineage depth replaces them on the DAG. (6)~Betweenness is exact
and directed. (7)~The within-family shuffle (A2) was added. (8)~RQ4 is framed as
a hypothetical flaw with full inheritance. (9)~The network-formation study was
added after the Topic~5-5 lecture.

% ---------------------------------------------------------------------------
\section{Step 5: Knowledge Discovered}

\subsection{RQ1: Merges join the network into one giant component}

Without merges the network is a forest of separate families: 60,177 weakly
connected components, the largest with 43,660 models. Adding the 48,030 merge
links reduces the count to 41,637 and grows the largest component 9.6 times, to
418,549 models (44.4\% of all nodes), which holds 2,448 roots from 28 named
families (Fig.~\ref{fig:components}). The next components are much smaller
(63,232, 25,428, 21,870 and 13,644). Without merges, each of the five largest
components is the tree of a single root: \model{FLUX.1-dev} (43,660 models),
\model{Qwen1.5-0.5B}, \model{Qwen1.5-1.8B}, \model{gemma-2b} and
\model{Qwen2.5-7B}. With merges, 12 components contain two or more named
families; without merges, none does. Merges are 5.2\% of the links, but they
are the links that connect families. This is the giant-component effect of
Topic~5-4: a small number of links between groups joins the groups.

\begin{figure}[t]
\centering
\includegraphics[width=\columnwidth]{rq1_component_sizes.png}
\caption{Component sizes with and without merge links. Merges grow the largest
component from 43,660 to 418,549 models.}
\label{fig:components}
\end{figure}

\begin{figure}[t]
\centering
\includegraphics[width=\columnwidth]{rq1_children_distribution.png}
\caption{Children per model. Left: log-binned distribution. Right: cumulative
distribution with the fitted slope ($\alpha=1.92$ for $10\le k\le 7{,}291$).}
\label{fig:children}
\end{figure}

\textbf{Hubs.} Only 87,535 models (9.3\%) have a child, and the mean number
of children is 0.99. The distribution of
children is heavy-tailed, with a cumulative-slope exponent $\alpha=1.92$ for
$10\le k\le 7{,}291$ (maximum-likelihood cross-check: 1.86). An exponent below 2
means the mean is dominated by the largest hubs \cite{barabasi2016}. The share
is extreme: the top 94 models (0.01\%) hold 42.1\% of all links, the top 0.1\%
hold 69.0\%, the top 1\% hold 84.8\% and the top 5\% hold 95.7\%. Even among the models that have a
child, the top 1\% (875 models) hold 68.3\% of the links (Fig.~\ref{fig:children}). The largest hub,
\model{FLUX.1-dev}, has 43,421 children. The relations differ: fine-tunes and
adapters have $\alpha\approx1.9$, quantizations 2.20 and merges 2.60, so merge
parents are the least concentrated: the largest merge parent has 412 merge
children, against 42,702 adapter children for the largest adapter parent.

\textbf{Depth.} The mean depth of non-root models is 1.62 in $G_{\text{full}}$ against 1.39 in
$G_{\text{nomerge}}$, and the maximum is 151 against 39. 6,917 models lie at
depth 10 or more, against 185 without merges. Most models are direct
children of a root (70.5\% at depth 1, against 73.3\% without merges), but the
share at depth 5 or more more than doubles, from 1.3\% to 2.9\%. Merges are on
average 6.4 steps from a root (median 3): they are built on top of other
merges, which gives long chains. A merge has 3.17 parents on average and up to
97.

\subsection{RQ2: Merges stay inside families and inside close relatives}

\textbf{Family mixing.} 85.5\% of merges combine parents from different roots,
so merges do join lineages. But they almost never join families. Among merges
whose parents' families are known, 45.5\% are same-family and 1.3\% are
cross-family; the remaining 53.1\% involve an unknown parent. Counting only
determined merges, 2.9\% are cross-family. The most common cross-family pairs
are Llama with Mistral (82 merges), Gemma with Qwen (12) and Gemma with Llama
(7). Llama and Mistral models use closely related architectures, which makes
their weights easier to combine; most other family pairs are almost never
merged. In the cocitation network the family
assortativity is $r=0.743$ with all labels and 0.988 among named families, with
modularity $Q=0.561$ and $0.678$ (Fig.~\ref{fig:family}). 80.6\% of the
cocitation weight joins two models with the same label.

\begin{figure}[t]
\centering
\includegraphics[width=\columnwidth]{rq2_family_matrix.png}
\caption{Cocitation weight between families. Almost all weight lies on the
diagonal: merges combine models of the same family.}
\label{fig:family}
\end{figure}

\textbf{Structure against $G(n,m)$.} A random graph with the same $n$ and $m$
has clustering $c/(n-1)=0.0009$; the cocitation network has 0.581, about 630
times more. Its giant component holds 61.7\% of the nodes against 99.99\% expected, and
its mean path length is 6.39 against 4.13 from $\ln n/\ln c$ (4.38 measured on
sampled $G(n,m)$ graphs). The network is clustered into tight groups that are
connected by few links, which is the opposite of a random graph.

\textbf{Inbreeding.} 32.6\% of merges have two parents within one generation of
each other (parent and child, or siblings), 45.9\% within three, and 48.7\%
share some ancestor (Fig.~\ref{fig:baseline}). The two types contribute about
equally (22.8\% ancestor type, 25.9\% common-ancestor type at any depth): a
model is often merged with its own parent. A few prolific authors could inflate
this, so we kept one merge per author and parent set (10,993 merges by 1,840
authors; the ten most active authors made 32.5\% of them). Inbreeding then rises to 35.1\% within one generation and 53.8\% at any
depth, so it is not caused by a few authors repeating themselves; the family
assortativity is also almost unchanged (0.734 against 0.743).

\begin{table}[t]
\caption{Share of inbred merges by depth $d$, split by type, against the two
shuffles (totals; mean of 100 runs).}
\label{tab:inbreeding}
\centering
\footnotesize
\setlength{\tabcolsep}{4pt}
\begin{tabular}{@{}lrrrrr@{}}
\toprule
 & \multicolumn{3}{c}{Observed} & \multicolumn{2}{c}{Shuffled total} \\
\cmidrule(lr){2-4}\cmidrule(l){5-6}
$d$ & Total & Ancestor & Common & A & A2 \\
\midrule
1   & 32.6\% & 19.1\% & 13.5\% & 5.3\%  & 16.3\% \\
2   & 42.0\% & 21.4\% & 20.6\% & 10.6\% & 25.4\% \\
3   & 45.9\% & 22.2\% & 23.7\% & 13.8\% & 29.9\% \\
any & 48.7\% & 22.8\% & 25.9\% & 16.6\% & 34.2\% \\
\bottomrule
\end{tabular}
\end{table}

\textbf{Which kind of relative?} Table~\ref{tab:inbreeding} splits the result
by type. The ancestor type is the stronger signal: within one generation,
19.1\% of merges combine a model with its own parent, against 2.6\% in shuffle A
and 8.9\% in A2. The common-ancestor type at any depth is 25.9\% against 20.7\%
in A2 ($z=12.7$), the smallest gap in the study. Distant cousins are common
inside a family anyway; what merges add is a strong preference for a model
and its direct ancestor.

\textbf{Against the baselines.} Table~\ref{tab:baseline} shows that no shuffle
comes near the observed values. Shuffle A accepted 30.9\% of swap attempts and
changed 94.3\% of the merge links; A2 accepted 7.7\% and changed 64.3\%. Across
the 100 runs of each, no run reached any observed value; the $z$-scores of the
measures in Table~\ref{tab:baseline} range from 30 to 1,584. The within-family shuffle A2 matters most: it keeps both
the degrees and the families, and still gives half the observed inbreeding
(16.3\% vs.\ 32.6\%). So merges are not only same-family; they choose close
relatives inside a family.

\begin{table}[t]
\caption{Observed values against degree-preserving shuffles
(mean of 100 runs; no run reached the observed value).}
\label{tab:baseline}
\centering
\footnotesize
\setlength{\tabcolsep}{3pt}
\begin{tabular}{@{}lrrr@{}}
\toprule
Measure & Observed & Shuffled (A) & In family (A2) \\
\midrule
Inbred within 1 & 32.6\% & 5.3\% & 16.3\% ($z{=}66$) \\
Inbred, any depth & 48.7\% & 16.6\% & 34.2\% ($z{=}30$) \\
Copies earlier merge & 48.8\% & 3.9\% & 12.1\% \\
Exact repeat & 34.4\% & 1.2\% & 4.9\% \\
Assortativity $r$ (all, unweighted) & 0.751 & 0.046 & -- \\
\midrule
 & & Rewired (B) & \\
\midrule
Assortativity $r$ (all, unweighted) & 0.751 & $-0.002$ & \\
Assortativity $r$ (named, unweighted) & 0.979 & $-0.004$ & \\
Clustering & 0.581 & 0.023 & \\
\bottomrule
\end{tabular}
\end{table}

\begin{figure}[t]
\centering
\includegraphics[width=\columnwidth]{baseline_inbreeding.png}
\caption{Share of inbred merges by depth, observed and in the two shuffles.
The within-family shuffle keeps degrees and families, yet stays far below the
observed values.}
\label{fig:baseline}
\end{figure}

\subsection{RQ3: Hubs are roots; bridges are merges}

\textbf{Hubs.} The models with the most children are base models released by
large labs: \model{FLUX.1-dev} (43,421), \model{Qwen1.5-0.5B} (32,534),
\model{Qwen1.5-1.8B} (30,627), \model{gemma-2b} (24,034) and
\model{distilbert-base-uncased} (13,263). PageRank agrees with the number of
children (Spearman 0.938 among models with a child; 13 of the top 20 are shared), and 16 of the top 20 by
PageRank are roots. PageRank adds models whose children have many descendants,
e.g.\ instruction-tuned bases that are fine-tuned again. The top 20 models hold
24.5\% of all PageRank, and the top 100 are mostly Qwen (37), Llama (15) and
Gemma (14) models.

\textbf{Bridges.} Betweenness tells a different story. Only 42,468 models
have positive betweenness, since a model needs both descendants and ancestors
to lie on a path; merges are 20.4\% of them. Merges are 1.6\% of all
nodes but carry 82.6\% of all directed betweenness, and 85\% of the top 20, 88\%
of the top 100 and 86.9\% of the top 1,000 models by betweenness are merges
(Fig.~\ref{fig:between}). The
top bridge is \model{L3-Umbral-Mind-RP-v3.0-8B}, a Llama merge with 18 parents.
Part of this is built in, since a merge lies on paths to several parents, but
the size of the effect shows that the paths between lineages run through a
small set of community merges and not through the base models. In the
cocitation network the bridges are base models that are merged across
families: \model{Mistral-Nemo-Instruct-2407} (cocited with 246 models from five
families), \model{Llama-3.1-8B}, \model{Qwen2.5-7B-Instruct} and
\model{Mistral-7B-v0.1}; 54 of the top 100 have neighbours in two or more
families. By family, the top 100 cocitation bridges are Llama (35), Qwen (22),
unknown (19) and Mistral (15) models.

\begin{figure}[t]
\centering
\includegraphics[width=\columnwidth]{rq3_lineage_betweenness.png}
\caption{Top 15 models by directed betweenness in the lineage graph. Twelve of
them are merges.}
\label{fig:between}
\end{figure}

\textbf{Communities.} Leiden finds 1,204 communities (range 1,202--1,206 over
20 runs), 82 of them with at least ten members (range 79--83). Their
modularity is $Q=0.890$ (range 0.888--0.890),
much higher than the 0.561 of the family labels, so the families are too coarse.
The communities are 96.2\% pure in family (range 95.5--96.4\%): Leiden splits
families into sub-lineages but almost never mixes them. It also places unknown
models: 46.0\% of the unknown models (range 44.3--53.2\%) sit in a community
whose labelled members share one family, so their family can be inferred from
the links.

\subsection{RQ4: A flaw in a base model reaches far, but a small audit covers it}

We assume a hypothetical flaw (e.g.\ a backdoor or a licence problem) in a root
model that is passed to every descendant. This is an upper bound, since a merge
or fine-tune may remove a flaw. The reach of a root is the number of models
found by BFS along reversed edges.

Merges raise the total number of root--descendant pairs from 862,654 to
1,135,867 (+31.7\%). 3,591 roots gain reach, and the number of roots that reach
1,000 or more models rises from 121 to 198. The gain is uneven:
\model{Mistral-Nemo-Instruct-2407} grows by 221.5\% (1,134 to 3,646),
\model{gemma-2-9b-it} by 84.5\%, \model{Llama-3.1-70B} by 78.1\%,
\model{Meta-Llama-3-8B-Instruct} by 60.8\% (3,688 to 5,932) and
\model{Mistral-7B-v0.1} by 26.6\% (9,119 to 11,548), while \model{FLUX.1-dev}
gains almost nothing (43,659 to 43,673), because its descendants are rarely
merged. The median root still reaches only two models, so the
exposure is concentrated in a few lineages. 30,147 models (3.4\% of non-roots)
descend from two or more roots, 6,467 from ten or more, and one from 128 roots;
279 of them inherit from two or more named families.

\begin{figure}[t]
\centering
\includegraphics[width=\columnwidth]{rq4_audit_coverage.png}
\caption{Share of all exposure covered by auditing the top-$k$ models, by
PageRank, children and reach, against random models.}
\label{fig:audit}
\end{figure}

\textbf{Audit coverage.} Topic~4 contrasts hub failure with random failure. We
ask how much of the total exposure is covered by auditing $k$ models
(Fig.~\ref{fig:audit}). The top model by PageRank covers 4.6\%, the top 10 cover
21.3\%, the top 100 cover 52.1\% and the top 1,000 cover 78.2\%. Ranking by
children gives almost the same result (78.1\% at 1,000), and ranking by reach is
worse at large $k$ (69.1\%) because high-reach models overlap: a merge and
its parents reach many of the same models. Without merges, the same reach
ranking covers 75.3\% at 1,000, so merges make greedy auditing by reach less
efficient. Random models
with a child cover 0.21\% at $k=100$ and 2.28\% at $k=1{,}000$; random models of
any kind cover 0.44\% at 1,000. The network is fragile to its hubs and robust to
random failure, exactly as a heavy-tailed network should be.

\subsection{Network formation (Topic 5-5)}

\textbf{Preferential attachment.} Replaying the links in creation order, the
attachment rate grows almost linearly with the number of children already held:
$R(k)\propto k^{\beta}$ with $\beta=1.01$ for $1\le k<4{,}096$
(Fig.~\ref{fig:growth}). This is the rich-get-richer rule of Price's model and
the Barab\'asi--Albert model \cite{price1976,barabasi1999}. Two observations
qualify it. First, models in the largest bin of $k$ receive only 0.26 of the fitted rate, so
attachment saturates at the top. Second, Price's model predicts
$\alpha = 2 + a/c \ge 2$ with $c=1.036$ parents per new model, while we measure
1.92, slightly below the bound. Age explains little: the Spearman correlation
between age and children is 0.046 (0.079 among models with a child), and the 754
models with 100 or more children gained a median of only 14.6\% of their
children in their first 90 days. So the early-mover advantage of Price's model
is weak; new strong base models overtake old ones. This points to a fitness
effect, i.e.\ some models are simply more useful, which the lectures mention as
the main extension of preferential attachment. Creation times before
March~2022 share one placeholder date on the Hub, so the age analysis uses the
dated models only.

\begin{figure*}[t]
\centering
\includegraphics[width=0.85\textwidth]{rq1_growth.png}
\caption{Network formation. Left: attachment rate $R(k)$ against the number of
children already held, with the fitted slope $\beta=1.01$ and the linear
reference. Right: median and 90th-percentile number of children by model age.}
\label{fig:growth}
\end{figure*}

\textbf{Vertex copying.} Merges copy each other. 48.8\% of merges have a parent
set with Jaccard similarity of at least 0.5 to an earlier merge, and 34.4\%
repeat an earlier parent set exactly (56.5\% share at least two parents with
some earlier merge); the shuffles give 3.9\% and 1.2\% (12.1\%
and 4.9\% within families). Most copying is an author repeating a recipe
(68.6\% of copies), but 15.3\% of all merges copy another author. 1,847 parent
sets are reused, covering 7,022 merges. The most copied recipe across authors,
\model{mistral-ft-optimized-1218} with \model{NeuralHermes-2.5-Mistral-7B}, was
used in 143 merges by 127 authors. The single largest recipe was used 223
times by one author, and the third, two image-generation checkpoints
(\model{noobai-XL-1.1} with \model{Illustrious-XL-v2.0}), was used 102 times by
three authors, so copying is not limited to language models. This is the copying mechanism of Topic~5-5,
and it explains both the clustering of the cocitation network and part of the
inbreeding.

% ---------------------------------------------------------------------------
\section{Step 6: Understanding}

\textbf{What merges do.} The results fit one picture. The Hub grows by
preferential attachment around a few base models from large labs, which gives
the heavy tail of RQ1 and the hubs of RQ3. Merges then add a second layer: they
are made by the community, they copy earlier recipes, and they combine close
relatives inside one family. Because each merge has several parents, this
second layer is what joins lineages into one giant component, what carries the
paths between lineages (82.6\% of the betweenness), and what adds a third to the
reach of a flaw. Yet merges almost never cross families, so the network is one
component made of family-pure communities, held together by few cross-family
links. The degree-preserving shuffles show that none of this follows from the
degrees alone.

\textbf{Implications.} For supply-chain risk, the audit result is the practical
one: checking the top 100 models by PageRank covers half of all exposure, which
is about 240 times more than checking 100 random models with a child. Merges should be
part of any audit, because some roots (e.g.\ \model{Mistral-Nemo-Instruct-2407})
reach 3.2 times more models once merges are counted. For
studies of merging, the high inbreeding means that many merges combine closely
related weights; comparisons of merge methods should control for parent
relatedness. Leiden communities can fill in missing family labels for about
half of the unknown models.

% ---------------------------------------------------------------------------
\section{Limitations and Future Scope}

\subsection{Limitations}

\textbf{Self-declared lineage.} All links come from what authors write in the
model card. A model that does not declare its parent becomes a root even if it
was derived, so every count in this report is a lower bound, and the true
network is likely more connected than the one we measure.

\textbf{Family labels.} The 41 patterns are a modelling choice. 23.5\% of merge
parents and 14.7\% of non-root models stay unknown, so family mixing is
measured on determined merges only. The labels are also sensitive: two patterns
(\texttt{bert} matching ``cybertron'' and \texttt{nemotron} labelling Llama
fine-tunes) were corrected after the first run, which changed the cross-family
count from 433 to 203.

\textbf{Full inheritance.} RQ4 assumes that every descendant inherits a flaw.
This is an upper bound: a fine-tune or a merge can weaken or remove a flaw, and
a quantized copy may keep it exactly. The audit ranking is still useful, but
the exposure numbers are not predictions.

\textbf{One snapshot.} Growth is replayed from creation times, not observed
over time. Deleted models are missing except as ghost parents, and models made
before March~2022 share one placeholder date, so the age analysis covers only
later models.

\textbf{Built-in effects.} Some patterns are partly caused by construction. A
merge lies on paths to all its parents, which raises its betweenness, and a
merge with many parents creates a clique in the cocitation network, which
raises clustering. The degree-preserving shuffles control for degrees but not
for these effects, so we report them as upper bounds on the true signal.

\textbf{Causes of growth.} The attachment rate cannot separate preferential
attachment from fitness, because a useful model both gains children and has
many children.

\textbf{Methods outside the lectures.} We used assortativity and
degree-preserving shuffles (from Newman \cite{newman2010}, as the proposal
promised), and simple statistics to report them: Spearman rank correlation,
Jaccard similarity and $z$-scores.

\subsection{Future Scope}

\textbf{Time-resolved network.} The source data set is updated daily. A series
of snapshots would let us observe growth directly, measure the attachment rate
without replaying creation times, and test a fitness model of growth
\cite{barabasi2016} against pure preferential attachment.

\textbf{Recovering hidden links.} Undeclared parents could be recovered from the
weights themselves, as proposed for the model atlas \cite{horwitz2025}. This
would reduce the number of false roots and unknown families, and show how much
the lower bounds underestimate.

\textbf{Partial inheritance.} RQ4 could weight each link by how much of a
parent survives: for example, full weight for a quantized copy and a smaller
weight for a fine-tune or for one parent of a merge. The audit question then
becomes a weighted reach problem on the same DAG.

\textbf{Node attributes.} Adding downloads, likes and licences to the nodes
would show whether the hubs by structure are also the most used models, and
whether merges combine parents with incompatible licences, a real case of the
supply-chain risk studied here.

\textbf{Other layers.} Models also link to the data sets they were trained on.
A bipartite model--data set network (Topic~2) would extend the lineage graph
and show whether merged models also share training data.

% ---------------------------------------------------------------------------
\section{Conclusion}

We built the lineage network of 942,088 Hugging Face models and studied what
merges do to it. Merges are 5.2\% of the links but join the network into one
giant component, carry most of the betweenness and add 31.7\% to the reach of a
hypothetical flaw. They combine models of the same family (assortativity 0.99
among named families) and often of the same close lineage (48.7\% inbred), and
they copy earlier recipes; random shuffles that keep the degrees reproduce none
of this. The network grows by near-linear preferential attachment, but age
matters little, so model fitness is the most likely driver. The code, the
cleaned data and all metric outputs are in the project archive, and a single
command reproduces every number in this report.

% ---------------------------------------------------------------------------
\begin{thebibliography}{10}

\bibitem{laufer2025}
B.~Laufer, H.~Oderinwale, and J.~Kleinberg, ``Anatomy of a machine learning
ecosystem: 2 million models on Hugging Face,'' 2025, \emph{arXiv:2508.06811}.

\bibitem{horwitz2025}
E.~Horwitz, N.~Kurer, J.~Kahana, L.~Amar, and Y.~Hoshen, ``We should chart an
atlas of all the world's models,'' 2025, \emph{arXiv:2503.10633}.

\bibitem{wei2026}
X.~Wei, L.~Kang, X.~Pan, and J.~Yang, ``Identifying disruptive models in the
open-source LLM community,'' 2026, \emph{arXiv:2604.11618}.

\bibitem{barabasi2016}
A.-L.~Barab\'asi, \emph{Network Science}. Cambridge, U.K.: Cambridge Univ.\
Press, 2016.

\bibitem{newman2010}
M.~E.~J.~Newman, \emph{Networks: An Introduction}. Oxford, U.K.: Oxford Univ.\
Press, 2010.

\bibitem{price1976}
D.~J.~de~Solla~Price, ``A general theory of bibliometric and other cumulative
advantage processes,'' \emph{J. Amer. Soc. Inf. Sci.}, vol.~27, no.~5,
pp.~292--306, Sep.\ 1976.

\bibitem{barabasi1999}
A.-L.~Barab\'asi and R.~Albert, ``Emergence of scaling in random networks,''
\emph{Science}, vol.~286, no.~5439, pp.~509--512, Oct.\ 1999.

\bibitem{traag2019}
V.~A.~Traag, L.~Waltman, and N.~J.~van~Eck, ``From Louvain to Leiden:
guaranteeing well-connected communities,'' \emph{Sci. Rep.}, vol.~9, no.~1, Art.~no.~5233,
Mar.\ 2019.

\bibitem{igraph}
G.~Cs\'ardi and T.~Nepusz, ``The igraph software package for complex network
research,'' \emph{InterJournal, Complex Systems}, no.~1695, 2006.

\end{thebibliography}

\end{document}
